\documentclass[11pt,a4paper]{article}
\usepackage[margin=1.5cm]{geometry}
\usepackage{mdframed}
\usepackage{parskip}
\usepackage{titlesec}
\usepackage{xcolor}
\usepackage{array}
\usepackage{booktabs}

\definecolor{accentblue}{RGB}{0,80,160}
\definecolor{noteframe}{RGB}{230,238,250}

\pagestyle{plain}
\setlength{\parindent}{0pt}

\newcommand{\notecard}[2]{%
  \begin{mdframed}[backgroundcolor=noteframe, linecolor=accentblue,
                   linewidth=1pt, roundcorner=4pt,
                   innertopmargin=6pt, innerbottommargin=8pt,
                   innerleftmargin=10pt, innerrightmargin=10pt,
                   skipabove=8pt, skipbelow=0pt]
    {\small\textcolor{accentblue}{\textbf{#1}}}\\[4pt]
    \small #2
  \end{mdframed}
}

\begin{document}

\begin{center}
  {\large\textbf{Presenter Notes}}\\[2pt]
  {\small Functional Forecasting of Turkish Electricity Load Curves --- Zeynep Bozoklu}
\end{center}
\vspace{4pt}

\notecard{Slide 1 --- Title}{
Introduce the thesis topic: applying Functional Data Analysis to forecast Turkey's daily electricity load curve. Core result: our model beats the official TEİAŞ forecast by 19\% with near-zero p-value. Walk through 14 slides in roughly 20 minutes.
}

\notecard{Slide 2 --- The Forecasting Object}{
Point to the overnight trough (around hour 4--5) and the broad daytime plateau (hours 9--20). The shape is structurally the same every day; what changes is the level and sharpness. Note how the official forecast (dashed) tracks the actual well on average but lies systematically below it --- we will quantify this underprediction on the next slide.
}

\notecard{Slide 3 --- Why Accuracy Matters}{
The official forecast always undershoots actual demand. The gap widened sharply in 2024. Our rolling-origin model reduces the average bias by 78\%. Just state the empirical fact --- do not speculate about why the official underpredicts. This table motivates the research question: is there systematic improvement to be had?
}

\notecard{Slide 4 --- From 24 Numbers to a Curve: The FDA Pipeline}{
Walk through the boxes left to right, then down and back. The core idea: instead of 24 separate forecasting models, we compress to 2 scores, forecast those, and reconstruct the full curve. The B-spline step also lets us compute the derivative of the curve, capturing ramp speed --- a feature we add later. This is the methodological contribution.
}

\notecard{Slide 5 --- Data}{
Data from EPIAS, Turkey's electricity market operator, via public API. Having both actual load and the official TEİAŞ day-ahead forecast is key: it gives us a directly comparable operational benchmark. The evaluation is deliberately aligned with the operational benchmark: same target days, same hourly load object, and no future information.
}

\notecard{Slide 6 --- FPCA: Two Numbers Capture the Curve}{
Top panel (FPC1): nearly flat at 0.13--0.28, meaning a high score uniformly lifts the whole curve --- the ``how much electricity today'' factor. Bottom panel (FPC2): positive at hours 0--5 and 18--23, negative daytime --- the overnight-vs-daytime contrast. Three coloured lines (2019--21, 2020--22, 2022--24) nearly overlap, which empirically justifies using a fixed FPCA basis rather than re-estimating it dynamically.
}

\notecard{Slide 7 --- Score Forecasting Model}{
Lag 7 captures the dominant weekly pattern (Monday resembles last Monday); lag 1 adjusts for recent drift. All features in $z_d$ are known before day $d$ starts --- no lookahead. Derivative features come from differentiating the B-spline curve, capturing how fast demand was changing yesterday. Key message to committee: holiday indicators are non-negotiable; the model is competitive only once they are included.
}

\notecard{Slide 8 --- Evaluation: Rolling-Origin Design}{
Stress the no-leakage principle: the model never sees the future. This mirrors real operational deployment. The Diebold-Mariano test is appropriate because we compare two forecasts for the same target over the same sample --- it accounts for forecast error correlation. Citing Harvey, Leybourne and Newbold (1997) for the small-sample correction is good practice if the committee asks.
}

\notecard{Slide 9 --- Main Results}{
Walk the table top to bottom. The seasonal na\"ive shows the floor --- the problem is hard. Pure VAR dynamics are worse than the official. The model crosses into positive territory only when holiday and load-shape features are added. Derivative features add another 2 points. Report the DM result explicitly: both the holiday/load-shape row and derivative row beat the official at $p<0.001$, with near-zero p-values ($p \approx 10^{-135}$ and $10^{-163}$). The main model (bold row) is the baseline contribution; the residual correction row previews the next extension.
}

\notecard{Slide 10 --- Where the Model Wins and Where It Loses}{
The ramp-hour dip at hours 6--9 is the key diagnostic. Demand changes very quickly during this window, so small timing errors become large hourly errors. We do not know what gives the official forecast its advantage at ramp hours --- do not speculate. This gap motivates the residual correction step on the next slide.
}

\notecard{Slide 11 --- Residual Correction: Exploiting Persistent Errors}{
The ACF panel for Hour 6 is the punchline: base model errors are strongly autocorrelated (ACF $\approx$ 0.4 at lag 1, persisting 14 days). The orange series after correction drops sharply toward zero. The correction is estimated on past data only --- no leakage. Key message: the residual step wins at the ramp hours (6--9) where the base model lost.
}

\notecard{Slide 12 --- Random Forest Benchmark}{
Give the Random Forest result proper weight. On the 2023 benchmark sample, FPCA-OLS beats the official forecast (MAE 910 vs.\ 1052, DM $p \approx 1.8\times10^{-32}$), but Direct RF is much stronger (MAE 718, +31.8\% vs.\ official, DM $p \approx 1.6\times10^{-167}$). Direct RF also beats FPCA-OLS at all 24 hours, with pairwise DM $p \approx 1.2\times10^{-85}$. The framing is not that FDA is invalid; it is that direct nonlinear hourly models are a strong benchmark, while FDA remains the interpretable functional model evaluated across the full 2023--2025 rolling sample.
}

\notecard{Slide 13 --- Thesis Road Map}{
The empirical core is complete. Open items: (1) a formal literature review situating this within FDA and electricity forecasting, (2) writing up the ramp-hour diagnostic evidence carefully, and (3) probabilistic extensions if time allows. Be direct with the committee about what is done vs planned.
}

\notecard{Slide 14 --- References}{
Keep this slide up while answering questions. If asked about methodology, cite Kokoszka \& Reimherr (2017). If asked about the DM test, cite Diebold \& Mariano (1995). If asked about FDA applied to energy forecasting, Hyndman \& Ullah (2007) is the closest precedent.
}

\end{document}
